% ============================================================================
% TDD --- F1tenth PoliMI car
% Source: car-interface design proposal (24 Oct 2026), docs/software/{car,docker,setup,git,github-access}.md,
%         hardware/bom/master_bom.xlsx, simulation plan 01.4, docs/scope.md (OI-10, OI-13)
% ============================================================================
\documentclass[11pt,a4paper]{article}

\input{preamble}
\input{doc}
\docbrand{polimi}

\docsetup{Technical Design Document (TDD)}{F1tenth PoliMI}{The car: stack, car interface and deployment}{0.01}{F1tenth PoliMI team}

\newcommand{\prop}{\target{Proposed}}

\begin{document}
\docfrontmatter{
0.01 & 2026-10-25 & F1tenth PoliMI team & Initial draft from the car-interface design proposal and the runbooks \\
}

% ============================================================================
\section{Introduction}
% ============================================================================
\subsection{Purpose}
This document defines how the car's software is built: the layers of the stack, the car interface the stack talks to, how the stack is started, where the car's measured numbers live, and how it is deployed on the car and on a developer's machine.

\subsection{Scope}
\begin{itemize}
  \item \textbf{In scope:} the stack on the car, the car interface (topics and frames), the car's hardware layer, the repositories, the Docker images, the development standards.
  \item \textbf{Out of scope:} the simulator's internals (the gym, the bridge, \code{polimi\_sim}). Here the simulator appears only as the second implementation of the car interface (\code{hw:=sim}). It gets its own documents, \code{RASD\_Gym} and \code{TDD\_Gym} \tbd.
\end{itemize}

\subsection{References}
\begin{itemize}
  \item \code{RASD\_F1tenth} --- requirements refined here.
  \item Car-interface design proposal, 24 Oct 2026 (the source of \cref{sec:arch,sec:interfaces}).
  \item Simulation plan 01.4, migration to Jazzy, the JAX gym and the simulator in our fork.
  \item Runbooks in \code{docs/software/}: \code{setup.md}, \code{docker.md}, \code{car.md}, \code{git.md}, \code{github-access.md}.
\end{itemize}

\subsection{Status Conventions}
\dec\ agreed \quad \prop\ proposed, waiting for the team's confirmation \quad \target{V2} later phase \quad \tbd\ open.

\begin{caveat}
\textbf{Today versus target.} The images run ROS~2 Humble today and the drivers come from \code{humble-polimi}. Jazzy is decided (DEC-01) and the migration is planned (plan 01.4), not done. This document describes the Jazzy target; where today differs, it says so.
\end{caveat}

% ============================================================================
\section{Architecture}\label{sec:arch}
% ============================================================================
\begin{takeaway}
\textbf{Key point.} One stack, two cars. The stack talks to the car only through a fixed set of ROS topics and TF frames, the \term{car interface}. On the car, \code{f1tenth\_system}'s drivers and \code{orbbec\_camera} serve it. In the simulator, our fork of \code{f1tenth\_gym\_ros} serves the same interface. The stack cannot tell which one it is running on, and neither fork depends on anything of ours. \prop
\end{takeaway}

\subsection{Component View}
Three layers. The logic and the common layer are one stack and run unchanged everywhere. The hardware layer is exactly one of two, picked by the launch argument \code{hw:=car} or \code{hw:=sim}.

The seam sits below the mux, at \code{/ackermann\_cmd}, the topic the VESC reads. The RASD requires joystick takeover, the deadman and AEB to work on the car and in the simulator (FR-SIM-002), so the mux and the safety nodes belong to the common layer, not to the hardware.

\begin{docfigure}
\begin{tikzpicture}
  \draw[zone=accent] (-2.2,4.9) rectangle (13.5,7.0);
  \node[zlabel=accent] at (-2.2,7.0) {logic};
  \draw[zone=accent] (-2.2,2.5) rectangle (13.5,4.6);
  \node[zlabel=accent] at (-2.2,4.6) {common};
  \node[svc] (slam) at (0,5.8) {\code{slam\_toolbox}};
  \node[svc] (amcl) at (3,5.8) {AMCL};
  \node[svc] (ekf)  at (6,5.8) {EKF};
  \node[svc] (lin)  at (9,5.8) {line detection};
  \node[svc] (pid)  at (12,5.8) {PID / Pure Pursuit};
  \node[svc] (aeb)  at (3,3.4) {AEB};
  \node[svc, text width=2.6cm] (mux) at (7,3.4) {\code{ackermann\_mux}\\priorities, locks};
  \draw[dashed, accent, thick] (-2.2,1.9) -- (13.5,1.9);
  \node[lbl, text=accent, anchor=west] at (9.2,1.9) {car interface};
  \node[ext, text width=5.2cm] (car) at (2.5,0.6) {\textbf{car} (\code{hw:=car})\\\code{f1tenth\_system} drivers, \code{orbbec\_camera}};
  \node[ext, text width=5.2cm] (sim) at (9.5,0.6) {\textbf{simulator} (\code{hw:=sim})\\gym, bridge, \code{polimi\_sim}};
  \draw[http] (pid.south) |- node[lbl, pos=0.75, above] {\code{/drive}} (mux.east);
  \draw[http] (aeb.east) -- node[lbl, above] {brake} (mux.west);
  \draw[http] (mux.south) -- (7,1.4) -| (car.north);
  \draw[http] (7,1.4) -| (sim.north);
  \node[lbl, anchor=west] at (7.1,2.3) {\code{/ackermann\_cmd}};
  \draw[aux] (-1.8,1.4) -- node[lbl, pos=0.5, right] {sensors, \code{/odom}, TF} (-1.8,5.8) -- (slam.west);
\end{tikzpicture}
\caption{Layers. Only \code{/ackermann\_cmd} goes down through the car interface; sensors, odometry and TF come up with the same names whichever hardware runs.}\label{fig:layers}
\end{docfigure}

\subsection{Components}
\begin{tabularx}{\textwidth}{@{}L{3.9cm} L{4.3cm} Y@{}}
\toprule
\hd{Component} & \hd{Responsibility} & \hd{Interfaces} \\
\midrule
\endhead
\multicolumn{3}{@{}l}{\textit{Logic}} \\
\code{slam\_toolbox} & Map the room (FR-MAP) & \code{/scan}, TF \arrowto\ \code{/map} \\
AMCL & Pose in the room map (FR-LOC) & \code{/scan}, \code{/odom}, \code{/map} \arrowto\ TF \code{map}\arrowto\code{odom} \\
EKF (\code{robot\_localization}) & Fused \code{odom} \tbd[OI-08] & \code{/odom}, IMU \arrowto\ TF \code{odom}\arrowto\code{base\_link} \\
\code{polimi\_lines} & Line detection (FR-LIN) & camera image, info \arrowto\ \code{nav\_msgs/Path} \\
\code{polimi\_control} & PID, Pure Pursuit (FR-PTH) & path, TF \arrowto\ \code{/drive} \\
\multicolumn{3}{@{}l}{\textit{Common}} \\
\code{ackermann\_mux} & Pick the highest-priority input still sending; locks (FR-SAF) & \code{/teleop}, \code{/drive}, \code{/joy} \arrowto\ \code{/ackermann\_cmd} \\
\code{polimi\_safety} & AEB (FR-SAF-004) & \code{/scan}, \code{/odom} \arrowto\ mux input \\
\code{polimi\_bringup} & The car interface (README), launches, our configs & \code{stack.launch.py hw:=car|sim} \\
\multicolumn{3}{@{}l}{\textit{Hardware, car}} \\
\code{vesc\_driver}, \code{vesc\_ackermann} & Motor, servo, odometry (FR-DRV-001) & \code{/ackermann\_cmd} \arrowto\ VESC; VESC \arrowto\ \code{/odom} \\
\code{joy}, \code{joy\_teleop} & Joystick (FR-DRV-002) & \code{/joy}, \code{/teleop} \\
\code{orbbec\_camera} & SL450 LiDAR, Gemini 335L camera and IMU (FR-DRV-003) & \code{/scan}, image, info, IMU \\
\multicolumn{3}{@{}l}{\textit{Hardware, simulator} (see \code{TDD\_Gym})} \\
fork of \code{f1tenth\_gym\_ros} & Serves the same interface & as the car \\
\bottomrule
\end{tabularx}

\subsection{Runtime View}
\subsubsection{Command path}
Autonomy publishes \code{/drive}. A person drives through \code{/teleop}, from the joystick on the car or the keyboard in the simulator. The mux picks by priority (joystick 100, autonomy 10, as in today's \code{mux.yaml}); the safety locks of OI-09 sit in the same mux. Below the seam, the car turns \code{/ackermann\_cmd} into VESC commands; the simulator turns it into gym actions through its actuation model.

\begin{docfigure}
\begin{tikzpicture}
  \node[ext] (joy) at (0,3) {\code{joy\_node}\\\code{joy\_teleop}};
  \node[svc] (ctl) at (0,1.5) {controller};
  \node[svc] (aeb) at (0,0) {AEB};
  \node[svc, minimum height=3.8cm, text width=2.5cm] (mux) at (5.6,1.5) {\code{ackermann\_mux}\\[4pt]teleop 100\\drive 10\\locks\\\tbd[OI-09]};
  \draw[http] (joy.east) -- node[lbl, above] {\code{/teleop}, \code{/joy}} (joy.east -| mux.west);
  \draw[http] (ctl.east) -- node[lbl, above] {\code{/drive}} (ctl.east -| mux.west);
  \draw[http] (aeb.east) -- node[lbl, above] {brake} (aeb.east -| mux.west);
  \draw[dashed, accent, thick] (10.4,-0.8) -- (10.4,4.1);
  \node[lbl, text=accent, anchor=south] at (10.4,4.1) {car interface};
  \node[ext, text width=3.0cm] (vesc) at (12.6,2.8) {\code{ackermann\_to\_vesc}\\VESC};
  \node[ext, text width=3.0cm] (act) at (12.6,0.2) {actuation model\\gym};
  \draw[http] (mux.east) -- node[lbl, above] {\code{/ackermann\_cmd}} (10.7,1.5) |- (vesc.west);
  \draw[http] (10.7,1.5) |- (act.west);
\end{tikzpicture}
\caption{Command path. Grey boxes are hardware-specific: in the simulator a keyboard bridge publishes \code{/teleop} in place of the joystick nodes, and the car's VESC is replaced by the actuation model. Everything else is identical on car and simulator, so takeover, deadman and AEB are tested in the simulator as they run on the car.}\label{fig:cmd}
\end{docfigure}

\subsubsection{Sensor path}
Each topic in the middle column is part of the car interface. The car and the simulator each have one producer per topic; the stack subscribes only to the middle column. The simulator's ground truth stays out of it: it is for debugging and scoring only, never read by the stack (FR-SIM-003).

\begin{docfigure}
\begin{tikzpicture}
  \foreach \y/\c/\t/\s [count=\i] in {
      4.4/{\code{orbbec\_camera} LiDAR}/{\code{/scan}, frame \code{laser}}/{LiDAR model},
      3.3/{\code{vesc\_to\_odom}}/{\code{/odom}, TF \code{odom}\arrowto\code{base\_link}}/{odometry model},
      2.2/{\code{orbbec\_camera} Gemini}/{camera image, info}/{camera model (01.5)},
      1.1/{\code{orbbec\_camera} Gemini}/{IMU}/{IMU model (01.5)},
      0.0/{\code{car.launch.py}}/{static TF mounts}/{\code{polimi\_sim}}} {
    \node[ext, minimum width=3.2cm, minimum height=0.8cm] (c\i) at (0,\y) {\c};
    \node[data, minimum width=3.6cm, minimum height=0.8cm] (t\i) at (6,\y) {\t};
    \node[ext, minimum width=3.2cm, minimum height=0.8cm] (s\i) at (12,\y) {\s};
    \draw[http] (c\i.east) -- (t\i.west);
    \draw[http] (s\i.west) -- (t\i.east);
  }
  \node[zlabel=darkgray, anchor=south] at (0,5.0) {car (\code{hw:=car})};
  \node[zlabel=accent, anchor=south] at (6,5.0) {car interface};
  \node[zlabel=darkgray, anchor=south] at (12,5.0) {simulator (\code{hw:=sim})};
\end{tikzpicture}
\caption{Sensor path. Same topic, same type, same frame, two producers. The mount TF comes from the same measured files on both sides (DEC-09), so the simulator places the LiDAR where the car's TF says it is.}\label{fig:sensors}
\end{docfigure}

\subsection{Deployment View}
\subsubsection{Machines and images}
One Docker image per machine type, built from \code{docker/ros/Dockerfile}. The repo's code, configs and maps are bind-mounted, so code is edited on the host and run in the container; \code{build/}, \code{install/} and \code{log/} live in Docker volumes.

\begin{tabularx}{\textwidth}{@{}L{1.4cm} L{3.2cm} L{1.8cm} Y@{}}
\toprule
\hd{Service} & \hd{Machine} & \hd{Simulator} & \hd{Networking} \\
\midrule
\code{arm} & Apple Silicon Mac & yes & port 8765 (Foxglove), bound to 127.0.0.1 \\
\code{x86} & Linux or Windows PC & yes & port 8765 (Foxglove), bound to 127.0.0.1 \\
\code{jet} & Jetson Orin Nano on the car & no & host \\
\code{novnc} & optional, RViz in the browser & -- & port 8080 \\
\bottomrule
\end{tabularx}

Base image \code{ros:jazzy} on Ubuntu 24.04 (DEC-01); today \code{ros:humble}. Middleware \code{rmw\_zenoh\_cpp}; the Zenoh router starts with the container.

\subsubsection{The car}
\begin{tabularx}{\textwidth}{@{}L{3.6cm} L{3.4cm} Y@{}}
\toprule
\hd{Device} & \hd{Link to the Orin Nano} & \hd{Driver} \\
\midrule
Orbbec Pulsar SL450 & Ethernet (the only port) & \code{orbbec\_camera}, \code{lidar.launch.py} \\
Orbbec Gemini 335L & USB 3 & \code{orbbec\_camera}, \code{gemini\_330\_series.launch.py} \\
VESC 6 MkVI & USB, \code{/dev/sensors/vesc} (udev rule \tbd) & \code{vesc\_driver} \\
8BitDo Ultimate 2C & wireless dongle, USB & \code{joy} \\
\bottomrule
\end{tabularx}

\subsubsection{Repositories}
\begin{tabularx}{\textwidth}{@{}L{3.7cm} L{4.9cm} Y@{}}
\toprule
\hd{Repo (\code{tfc-org})} & \hd{Submodule path} & \hd{Holds} \\
\midrule
\code{f1tenth\_polimi} & -- (this repo) & Our packages, configs, maps, Docker, docs \\
\code{f1tenth\_system\_polimi} & \code{ros2\_ws/src/f1tenth\_system} & Upstream drivers, branch \code{jazzy-polimi} (today \code{humble-polimi}) \\
\code{f1tenth\_gym\_polimi} & \code{ros2\_ws/src/f1tenth\_gym\_ros} (plan 01.4) & Upstream bridge plus \code{polimi\_sim}, branch \code{dev-jazzy-polimi} \\
\bottomrule
\end{tabularx}

Dependencies point one way only: from our repo outwards. Neither fork knows our stack exists (DEC-08).

\subsubsection{Launch}
One entry point, \code{stack.launch.py}, with \code{hw:=car} or \code{hw:=sim}. It always starts the common layer and the logic asked for, plus exactly one hardware launch. \prop
\begin{itemize}
  \item \code{stack.launch.py hw:=car|sim}
  \begin{itemize}
    \item \code{common.launch.py}: \code{ackermann\_mux} with our \code{mux.yaml}, safety.
    \item \code{car.launch.py} (\code{hw:=car}): \code{f1tenth\_system}'s driver nodes one by one, \code{orbbec\_camera}, the mount TF from \code{config/car/}.
    \item \code{sim.launch.py} (\code{hw:=sim}): includes the fork's \code{sim\_car.launch.py} with \code{teleop\_topic:=teleop} and \code{config\_dir:=\textasciitilde/ws/config}, nothing else.
    \item the logic: SLAM or localization, line detection, controller, as asked.
  \end{itemize}
\end{itemize}
Today the car starts with \code{f1tenth\_stack}'s \code{no\_lidar\_bringup\_launch.py} and the simulator with \code{simcar}, which becomes a shortcut for \code{hw:=sim}.

\subsection{Design Decisions}
\begin{tabularx}{\textwidth}{@{}L{1.4cm} L{4.2cm} Y@{}}
\toprule
\hd{ID} & \hd{Decision} & \hd{Rationale / alternatives} \\
\midrule
\endhead
DEC-01 & ROS~2 Jazzy on every image \dec & Humble ends on 31 May 2027, before or during the summer school; Jazzy is supported to 31 May 2029. Lyrical is too new: no \code{f1tenth\_gym\_ros} branch, no JetPack for its Ubuntu. (OI-13, 21 Oct 2026.) \\
DEC-02 & One image rebuild adds \code{nav2\_amcl}, \code{robot\_localization}, \code{image\_transport\_plugins}, \code{cv\_bridge} \dec & In the Jazzy base apt layer, so AMCL and the EKF can be tried before any stack package exists. Python nodes never import \code{cv\_bridge} (it breaks under the gym's NumPy~2) and convert images with NumPy. (OI-10, 22 Oct 2026.) \\
DEC-03 & Seam below the mux, at \code{/ackermann\_cmd} \prop & The simulator replaces the VESC, the motor, the chassis and the sensors, nothing above them. With the mux above the seam, takeover, deadman and AEB are the same nodes in both places. (R1) \\
DEC-04 & Plain topics are the contract; no extension mechanism \prop & Our launch includes the fork's launch with two arguments; one switch, \code{hw:=}. No new code in the fork, and the contract can be checked with \code{ros2 topic info} on both sides. (R2) \\
DEC-05 & \code{polimi\_bringup} holds the contract and the launches; one package per function for the logic \prop & Functions have different dependencies (the camera stack versus none) and their own tests; one bringup package keeps how it all starts in one place. (R3) \\
DEC-06 & Our own \code{car.launch.py}, not \code{f1tenth\_stack}'s bringup \prop & Upstream's \code{bringup\_launch.py} starts the Hokuyo driver, its own mux and a hard-coded laser TF, none of which match our car. \code{f1tenth\_system} stays a pure upstream fork: the scope's ``our fork, extended with all our drivers'' becomes ``its driver packages, launched by our bringup''. (R4) \\
DEC-07 & Our own \code{mux.yaml}, in \code{polimi\_bringup/config/} \prop & Starts from today's values (joystick 100, autonomy 10). OI-09 adds the e-stop and deadman locks and the AEB input, and those are ours, not upstream's. (R5) \\
DEC-08 & The simulator fork runs from a bare clone \prop & \code{polimi\_sim} drops the mux node and its \code{ackermann\_mux} and \code{f1tenth\_stack} dependencies, and gains \code{teleop\_topic} (default \code{ackermann\_cmd}). Replaces decision~8 of plan 01.4, which made the fork depend on \code{f1tenth\_system}. (R6) \\
DEC-09 & The hardware side publishes the mount TF, from the same files \prop & Car: \code{car.launch.py} from \code{config/car/*.yaml}. Simulator: \code{polimi\_sim} from the same files via \code{config\_dir}. One file, two readers. (R7) \\
DEC-10 & \code{odom}\arrowto\code{base\_link} has one owner at a time \prop & The odometry publishes it until the EKF lands (OI-08); then the EKF does, and \code{publish\_tf} goes off on car and simulator in the same change. Two publishers make the tree jump. (R8) \\
DEC-11 & Everything runs in Docker, code bind-mounted \dec & Same images on Mac, PC and car; the host needs only Docker and git. \\
DEC-12 & Forks stay upstream plus upstreamable fixes; our code in \code{polimi\_*} packages \dec & Upstream fixes can be pulled; only fixes we cannot make from outside go into a fork. \\
\bottomrule
\end{tabularx}

% ============================================================================
\section{Data}
% ============================================================================
The measured car lives in \code{ros2\_ws/config/}, written once and read by both the car and the simulator. The simulator fork's defaults only fill gaps and never hold our measurements; it logs where each value came from.

\begin{tabularx}{\textwidth}{@{}L{5.6cm} L{2.6cm} Y@{}}
\toprule
\hd{File} & \hd{Readers} & \hd{Holds} \\
\midrule
\endhead
\code{config/car/laser\_mount.yaml} & car bringup, simulator & \code{base\_link}\arrowto\code{laser}: today f1tenth's defaults 0.27, 0, 0.11\,m, [MEASURE] on the car \\
\code{config/car/camera\_mount.yaml} \tbd & car bringup, simulator & \code{base\_link}\arrowto\code{camera\_link} \\
\code{config/lidar/sl450.yaml} & simulator & SL450 model: 40\,Hz, 270\textdegree, 0.2\textdegree, ranges 0.05--25\,m, noise, dropout \\
\code{config/maps/*.yaml} & simulator & track maps for the simulator \\
\code{maps/} & SLAM, AMCL & saved room maps \\
\code{polimi\_bringup/config/mux.yaml} \prop & mux & priorities, timeouts (0.2\,s), locks \\
\code{polimi\_bringup/config/} slam, amcl, ekf \prop & logic & node parameters \\
\code{vesc.yaml} (\code{f1tenth\_stack}) & VESC driver & gains, servo offset, wheelbase \\
\bottomrule
\end{tabularx}

% ============================================================================
\section{Interfaces}\label{sec:interfaces}
% ============================================================================
\subsection{The car interface}
This table is the contract the stack is built against; it becomes the README of \code{polimi\_bringup} \prop. All message types are standard ROS ones, so no interface package is needed. Exact camera and IMU topic names follow \code{orbbec\_camera}'s defaults, so the car needs no remaps \tbd[T-02].

Only \code{/ackermann\_cmd} goes from the stack to the hardware; every other row comes from the hardware.

\begin{tabularx}{\textwidth}{@{}L{3.6cm} L{3.7cm} Y Y@{}}
\toprule
\hd{Topic / frame} & \hd{Type} & \hd{Car} & \hd{Simulator} \\
\midrule
\endhead
\code{/ackermann\_cmd} & \code{AckermannDriveStamped} & \code{ackermann\_to\_vesc} & actuation model \\
\code{/teleop} & \code{AckermannDriveStamped} & \code{joy\_teleop} & keyboard bridge \\
\code{/joy} & \code{Joy} & \code{joy\_node} & stand-in \tbd[T-01] \\
\code{/scan}, frame \code{laser} & \code{LaserScan} & \code{orbbec\_camera} LiDAR & LiDAR model \\
\code{/odom} & \code{Odometry} & \code{vesc\_to\_odom} & odometry model \\
TF \code{odom}\arrowto\code{base\_link} & TF & \code{vesc\_to\_odom} & odometry model \\
camera image, info & \code{Image}, \code{CameraInfo} & \code{orbbec\_camera} & camera model (plan 01.5) \\
IMU & \code{Imu} & \code{orbbec\_camera} & IMU model (plan 01.5) \\
TF \code{base\_link}\arrowto\code{laser}, \code{camera\_link} & static TF & \code{car.launch.py} & \code{polimi\_sim} \\
\bottomrule
\end{tabularx}

\code{/drive} (\code{AckermannDriveStamped}) is inside the stack, not part of the car interface: any controller writes it, and it enters the mux at priority 10.

\subsection{Path following}
\begin{itemize}
  \item \textbf{Input path:} \code{nav\_msgs/Path} on one topic, any frame; offline paths latched, online paths recomputed per camera frame (FR-PTH-002, 005).
  \item \textbf{Input pose:} TF lookup from the path's frame to \code{base\_link} (FR-PTH-003).
  \item \textbf{Output:} \code{/drive}. A stale online path stops the car (FR-PTH-007).
\end{itemize}

% ============================================================================
\section{Security \& Privacy}
% ============================================================================
\begin{itemize}
  \item \textbf{Physical safety comes first:} the mux, its locks and AEB (FR-SAF) are the safety design; they run on the car and in the simulator unchanged (DEC-03).
  \item \textbf{Network:} on Mac and PC the container's ports bind to 127.0.0.1 only. The car's container uses host networking; ROS traffic on the car's network is not authenticated \tbd.
  \item \textbf{Secrets and personal data:} none in the stack. Camera images may show people in the room; nothing is recorded by default.
  \item \textbf{Repository access:} protected branches and required reviews, as in the runbook \code{github-access.md}.
\end{itemize}

% ============================================================================
\section{Requirements Traceability}
% ============================================================================
\begin{tabularx}{\textwidth}{@{}L{2.6cm} L{4.4cm} Y@{}}
\toprule
\hd{Requirement} & \hd{Component(s)} & \hd{Test(s)} \\
\midrule
\endhead
\code{FR-DRV-001} & \code{vesc\_driver}, \code{vesc\_ackermann} & bench: \code{ros2 topic echo /odom} while driving by joystick \\
\code{FR-DRV-002} & \code{joy}, \code{joy\_teleop} & bench: \code{/teleop} only while the enable button is held \\
\code{FR-DRV-003} & \code{orbbec\_camera} & bench: rates of \code{/scan}, image, IMU \\
\code{FR-DRV-004} & \code{car.launch.py} (DEC-09) & \code{tf2\_echo base\_link laser} matches \code{laser\_mount.yaml} \\
\code{FR-MAP-001}, \code{002} & \code{slam\_toolbox} & AC-V1-01 \\
\code{FR-LOC-001}, \code{002} & AMCL & AC-V1-01 \\
\code{FR-LOC-003} & EKF (DEC-10) & \tbd[OI-08] \\
\code{FR-LIN-001}, \code{002} & \code{polimi\_lines} & recorded camera data; AC-V1-02 \\
\code{FR-PTH-001} to \code{008} & \code{polimi\_control} & unit tests on the error and window search; AC-V1-02, 03 \\
\code{FR-SAF-001} to \code{005} & \code{ackermann\_mux}, \code{mux.yaml}, \code{polimi\_safety} & AC-V1-04, in the simulator first \\
\code{FR-SIM-001} to \code{003} & car interface, \code{hw:=sim} & \code{ros2 topic info} on both sides; AC-V1-02 \\
\code{FR-DOC-001} & runbooks, RASD, TDD & review by someone outside the team \\
\code{FR-ADV-001} to \code{003} & \target{V2} \tbd[OI-01] & AC-V2-A \\
\code{FR-ADV-011} to \code{013} & \target{V2} \tbd[OI-01] & AC-V2-B \\
\code{NFR-POR-001} & images, DEC-01, DEC-11 & the same launch on \code{arm}, \code{x86}, \code{jet} \\
\code{NFR-VAR-001}, \code{002} & logic & laps on more than one layout and room \\
\code{NFR-VAR-003} & \code{polimi\_lines} & \tbd[OI-06], once floor and tape are known \\
\code{NFR-PER-001} & logic, controller & \tbd: no target yet \\
\code{NFR-PER-002} & \code{orbbec\_camera} LiDAR & bench: \code{ros2 topic hz /scan} \\
\code{NFR-PER-003} & \code{orbbec\_camera} camera & bench: \code{ros2 topic hz} on the image \\
\code{NFR-PER-004} & \code{orbbec\_camera}, \code{vesc\_driver} & \tbd: method to measure timestamp error \\
\code{NFR-PER-005} & \code{mux.yaml} & stop sending \code{/drive}: the car stops within 0.2\,s \\
\code{NFR-SAF-001} & mux, locks, AEB & AC-V1-04, in the simulator and on the car \\
\code{NFR-SAF-002} & test plan & AC-V1-04 passed on the car before autonomous laps \\
\code{NFR-MNT-001}, \code{002} & development standards (DEC-12) & PR review \\
\code{FR-OFF-001} & \tbd[OI-11] & optional \\
\bottomrule
\end{tabularx}

% ============================================================================
\section{Implementation}
% ============================================================================
\subsection{Technology Stack}
\begin{tabularx}{\textwidth}{@{}L{3cm} L{4.4cm} Y@{}}
\toprule
\hd{Layer} & \hd{Choice} & \hd{Why} \\
\midrule
\endhead
Compute & Jetson Orin Nano Super, 500\,GB NVMe & BOM; fits the F1TENTH deck \\
OS, container & Ubuntu 24.04 images, Docker; JetPack 6 on the car \tbd[T-10] & DEC-01, DEC-11 \\
Middleware & ROS~2 Jazzy, \code{rmw\_zenoh\_cpp} & DEC-01 \\
Drivers & \code{f1tenth\_system} (\code{vesc}, \code{ackermann\_mux}, \code{teleop\_tools}), \code{orbbec\_camera} & C-022 \\
Mapping, localization & \code{slam\_toolbox}, \code{nav2\_amcl}, \code{robot\_localization} & scope; DEC-02 \\
Visualization & Foxglove bridge; RViz through noVNC & runs without a display on the host \\
Python & system Python of the image; \code{uv} on the host & house rule \\
\bottomrule
\end{tabularx}

\subsection{Plan}
\begin{enumerate}
  \item \textbf{Plan 01.4:} Jazzy images, the JAX gym, the simulator fork as a submodule.
  \item \textbf{\code{polimi\_bringup}:} the car interface README, \code{stack.launch.py}, \code{sim.launch.py}, \code{common.launch.py}, our \code{mux.yaml}; \code{simcar} calls \code{hw:=sim}.
  \item \textbf{Safety:} OI-09 locks and AEB, tested in the simulator.
  \item \textbf{Logic:} mapping and localization, line detection, path following, as the base platform needs.
  \item \textbf{Car bringup}, when the car is built: \code{car.launch.py}, udev rules, measured configs, safety tests on the car.
  \item \textbf{Advanced topic} \target{V2}, once OI-01 is decided.
\end{enumerate}

\subsection{Development Standards}
\begin{itemize}
  \item \textbf{Branches and reviews:} every change through a PR. Protected branches (\code{main}; the forks' \code{humble-*}, \code{jazzy-*}, \code{dev-humble*}, \code{dev-jazzy*}) block force pushes and deletion and need one approval.
  \item \textbf{Submodules:} changing a fork takes a commit there, then a pointer commit here; \code{push.recurseSubmodules=check} blocks a pointer to an unpushed commit. \code{scripts/setup\_git.sh} sets this up and adds the \code{upstream} remotes.
  \item \textbf{Upstream sync:} merge upstream into our branch, never rewrite it; preview with \code{git log ours..upstream}.
  \item \textbf{Pins:} gym and bridge refs are bumped on purpose, never by following a branch (NFR-MNT-002).
  \item \textbf{Rebuilds:} code and configs need only \code{cb}; a \code{package.xml}, the Dockerfile or a pin needs an image rebuild.
  \item \textbf{Python:} \code{uv} on the host, never \code{pip install} directly.
\end{itemize}

% ============================================================================
\section{Testing}
% ============================================================================
\subsection{Strategy}
\begin{tabularx}{\textwidth}{@{}L{3cm} L{4cm} Y@{}}
\toprule
\hd{Level} & \hd{Scope} & \hd{Tooling / gate} \\
\midrule
Unit & Controller math, line fitting, AEB time to collision & \code{colcon test}, pytest \\
Contract & The car interface on both sides & \code{ros2 topic info} and \code{tf2\_echo}: same names, types, frames on \code{hw:=car} and \code{hw:=sim} \\
System, simulator & Full flows: map, localize, laps, safety & AC-V1-01 to 04 in the simulator first \\
System, car & The same flows on the car & AC-V1-01 to 04 on the car; safety before autonomous laps (NFR-SAF-002) \\
\bottomrule
\end{tabularx}

\subsection{Acceptance Criteria}
As in \code{RASD\_F1tenth}: AC-V1-01 to 04 for the base platform, AC-V2-A or AC-V2-B for the advanced topic. Leaving the lines is judged by the people watching on the car, and measured in the simulator.

% ============================================================================
\section{Open Items}
% ============================================================================
\begin{tabularx}{\textwidth}{@{}L{1.2cm} L{4.2cm} Y@{}}
\toprule
\hd{ID} & \hd{Topic} & \hd{Notes} \\
\midrule
\endhead
T-01 & Deadman and e-stop in the simulator & Docker Desktop on macOS has no USB passthrough, so there is no joystick in the container. The locks need a stand-in for \code{/joy}: a keyboard key that toggles the deadman, or a joystick read on the Mac and bridged in. Settle with OI-09. \\
T-02 & Camera and IMU topic names [VERIFY] & From \code{orbbec\_camera}'s Gemini launch, once the camera is on the bench; the simulator's camera model uses the same names. \\
T-03 & LiDAR frame name [VERIFY] & From \code{orbbec\_camera}'s \code{lidar.launch.py}; set it to \code{laser} in \code{car.launch.py} if it differs. \\
T-04 & Confirm DEC-03 to DEC-10 & The proposal amends plan 01.4: decision 8 withdrawn, M2 and M3 change, the scope's ``Drivers'' line is reworded. \\
T-05 & Wheelbase & \code{vesc.yaml} has 0.25\,m, the simulated odometry the gym's 0.3302\,m. Measure the car and align both. \\
T-06 & Mount measurements [MEASURE] & \code{laser\_mount.yaml} holds f1tenth's defaults; the camera mount file does not exist yet. \\
T-07 & VESC device & udev rule on the Jetson for \code{/dev/sensors/vesc}, expected by \code{vesc.yaml}. \\
T-08 & \code{orbbec\_camera} version & 2.9.3 is in the Jazzy apt repo; check it with the SL450's firmware (plan 01.4, M-O8). \\
T-10 & JetPack and Ubuntu 24.04 [VERIFY] & JetPack 6 is based on Ubuntu 22.04. Check that the \code{jet} service runs a \code{ros:jazzy} (24.04) container with \code{runtime: nvidia} on the Orin Nano. \\
T-09 & Unmatched remap & Today's bringup remaps \code{ackermann\_cmd\_out}\arrowto\code{ackermann\_drive}, which matches nothing: the mux publishes \code{ackermann\_cmd}. Gone with DEC-06. \\
\bottomrule
\end{tabularx}

\end{document}
